% --- Archon physics formalization source begin ---
% archon:physics
% archon:covers PhyXMiniProblems/problem_phyx_mini_0566.lean
% archon:source-report reports/phyx_mini/problem_phyx_mini_0566.source.json

\chapter{Physics problem phyx\_mini\_0566}
\label{ch:PhyXMiniProblems_problem_phyx_mini_0566}

\paragraph{Problem source.}
\#\# Physical scenario

Observer O measures the area 3 m\$\textasciicircum{}2\$ of an equilateral triangle ABC at rest in his xy-plane. An observer O′ moving relative to O at 0.6c observes the triangle A′B′C′. The triangle symmetry axis forms an angle \$\textbackslash{}theta = 15\textasciicircum{}\{\textbackslash{}circ\}\$ with respect to the direction x–x′ as measured by observer O

\#\# Question

What is the area measured by O′?

\#\# Answer choices

- A: A: 3 m\$\textasciicircum{}2\$

- B: B: 1.4 m\$\textasciicircum{}2\$

- C: C: 3.3 m\$\textasciicircum{}2\$

- D: D: 3.2 m\$\textasciicircum{}2\$

\#\# Figure caption (auxiliary; use the image as primary evidence)

The image depicts two coordinate systems: one on the left and one on the right. Both systems display red triangles.

\#\#\# Left Side:
- **Axes**: The axes are labeled as \textbackslash{}( x \textbackslash{}) (horizontal) and \textbackslash{}( y \textbackslash{}) (vertical).
- **Triangle**:
  - Points are labeled as \textbackslash{}( A \textbackslash{}), \textbackslash{}( B \textbackslash{}), and \textbackslash{}( C \textbackslash{}), forming a triangle.
  - An angle \textbackslash{}( \textbackslash{}Theta \textbackslash{}) is shown between point \textbackslash{}( C \textbackslash{}) and point \textbackslash{}( A \textbackslash{}).

\#\#\# Right Side:
- **Axes**: The axes are labeled as \textbackslash{}( x' \textbackslash{}) (horizontal) and \textbackslash{}( y' \textbackslash{}) (vertical).
- **Vector and Velocity**:
  - There is a vector labeled \textbackslash{}( V \textbackslash{}) along the \textbackslash{}( x' \textbackslash{}) axis.
- **Triangle**:
  - Points are labeled as \textbackslash{}( A' \textbackslash{}), \textbackslash{}( B' \textbackslash{}), and \textbackslash{}( C' \textbackslash{}).

\#\#\# Additional Details:
- The \textbackslash{}( x \textbackslash{}) axis is noted as being equivalent to the \textbackslash{}( x' \textbackslash{}) axis (\textbackslash{}( x \textbackslash{}equiv x' \textbackslash{})).
- The triangles in both coordinate systems are similar in shape but oriented differently. 

This setup suggests a transformation or rotation between the two coordinate systems.

\#\# Dataset metadata

Relativity; Physical Model Grounding Reasoning, Spatial Relation Reasoning

\paragraph{Recorded answer/context.}
B: B: 1.4 m\$\textasciicircum{}2\$

\paragraph{Figure/image path.}
/root/proposal\_for\_physic/science-mango/phyx\_mini\_run/phyx\_data/test\_image/566.png

\paragraph{Formalization target.}
create a compiling Lean file with sorry bodies at `PhyXMiniProblems/problem\_phyx\_mini\_0566.lean`. The Lean declarations must preserve the physical quantities, dimensions or dimensional roles, figure labels, governing-law hypotheses, and final relation expressed by this problem.
Use Mathlib/Physlib names found through LeanExplore where available. If a physics API is missing, introduce faithful local abstractions rather than scalar placeholder aliases.

\begin{theorem}[Physics formalization target]
\label{thm:physics:phyx_mini_0566:target}
\lean{PhyXMiniProblems.ProblemPhyXMini0566.problem_phyx_mini_0566}
\uses{def:physics:phyx-mini-0566:phyxminiproblems-problemphyxmini0566-areainsquaremeters, def:physics:phyx-mini-0566:phyxminiproblems-problemphyxmini0566-relativistictrianglesetup, def:physics:phyx-mini-0566:phyxminiproblems-problemphyxmini0566-matchesrelativistictrianglescenario, def:physics:phyx-mini-0566:phyxminiproblems-problemphyxmini0566-matchessuppliedtrianglefigure, def:physics:phyx-mini-0566:phyxminiproblems-problemphyxmini0566-matchesproblemdata, def:physics:phyx-mini-0566:phyxminiproblems-problemphyxmini0566-satisfiesrestequilateralgeometry, def:physics:phyx-mini-0566:phyxminiproblems-problemphyxmini0566-hasphysicalrelativisticparameters, def:physics:phyx-mini-0566:phyxminiproblems-problemphyxmini0566-satisfiessimultaneouscoordinateareameasurement, def:physics:phyx-mini-0566:phyxminiproblems-problemphyxmini0566-obeyslorentzcoordinatecontraction}
The assigned autoformalize agent should translate this physics problem into Lean declarations in the covered file, with theorem and lemma proof bodies written as `by sorry`.
\end{theorem}
\begin{proof}
This is an autoformalization task, not a proof task. Produce statements that can later be proved by the physics prover without weakening the physical model.
\end{proof}

% --- Archon declaration topology begin ---
\section{Declaration topology}
\begin{definition}[Length Quantity]
  \label{def:physics:phyx-mini-0566:phyxminiproblems-problemphyxmini0566-lengthquantity}
  \lean{PhyXMiniProblems.ProblemPhyXMini0566.LengthQuantity}
  A nonnegative, unit-independent physical length
\end{definition}
\begin{proof}
  This is the declared model definition of Length Quantity.
\end{proof}
\begin{definition}[Area Quantity]
  \label{def:physics:phyx-mini-0566:phyxminiproblems-problemphyxmini0566-areaquantity}
  \lean{PhyXMiniProblems.ProblemPhyXMini0566.AreaQuantity}
  A nonnegative, unit-independent physical area
\end{definition}
\begin{proof}
  This is the declared model definition of Area Quantity.
\end{proof}
\begin{definition}[length In Meters]
  \label{def:physics:phyx-mini-0566:phyxminiproblems-problemphyxmini0566-lengthinmeters}
  \lean{PhyXMiniProblems.ProblemPhyXMini0566.lengthInMeters}
  \uses{def:physics:phyx-mini-0566:phyxminiproblems-problemphyxmini0566-lengthquantity}
  Read a physical length in coherent SI metres
\end{definition}
\begin{proof}
  This is the declared model definition of length In Meters.
\end{proof}
\begin{definition}[area In Square Meters]
  \label{def:physics:phyx-mini-0566:phyxminiproblems-problemphyxmini0566-areainsquaremeters}
  \lean{PhyXMiniProblems.ProblemPhyXMini0566.areaInSquareMeters}
  \uses{def:physics:phyx-mini-0566:phyxminiproblems-problemphyxmini0566-areaquantity}
  Read a physical area in coherent SI square metres
\end{definition}
\begin{proof}
  This is the declared model definition of area In Square Meters.
\end{proof}
\begin{definition}[speed In Meters Per Second]
  \label{def:physics:phyx-mini-0566:phyxminiproblems-problemphyxmini0566-speedinmeterspersecond}
  \lean{PhyXMiniProblems.ProblemPhyXMini0566.speedInMetersPerSecond}
  Read a physical speed in coherent SI metres per second
\end{definition}
\begin{proof}
  This is the declared model definition of speed In Meters Per Second.
\end{proof}
\begin{definition}[vacuum Speed Of Light In Meters Per Second]
  \label{def:physics:phyx-mini-0566:phyxminiproblems-problemphyxmini0566-vacuumspeedoflightinmeterspersecond}
  \lean{PhyXMiniProblems.ProblemPhyXMini0566.vacuumSpeedOfLightInMetersPerSecond}
  Physlib's exact vacuum speed of light, read in metres per second
\end{definition}
\begin{proof}
  This is the declared model definition of vacuum Speed Of Light In Meters Per Second.
\end{proof}
\begin{definition}[degrees To Radians]
  \label{def:physics:phyx-mini-0566:phyxminiproblems-problemphyxmini0566-degreestoradians}
  \lean{PhyXMiniProblems.ProblemPhyXMini0566.degreesToRadians}
  Convert a numerical degree readout to a radian coordinate
\end{definition}
\begin{proof}
  This is the declared model definition of degrees To Radians.
\end{proof}
\begin{definition}[Observer Label]
  \label{def:physics:phyx-mini-0566:phyxminiproblems-problemphyxmini0566-observerlabel}
  \lean{PhyXMiniProblems.ProblemPhyXMini0566.ObserverLabel}
  The rest observer and the relatively moving observer in the problem
\end{definition}
\begin{proof}
  This is the declared model definition of Observer Label.
\end{proof}
\begin{definition}[Triangle Vertex]
  \label{def:physics:phyx-mini-0566:phyxminiproblems-problemphyxmini0566-trianglevertex}
  \lean{PhyXMiniProblems.ProblemPhyXMini0566.TriangleVertex}
  A physical vertex, labelled without or with primes according to frame
\end{definition}
\begin{proof}
  This is the declared model definition of Triangle Vertex.
\end{proof}
\begin{definition}[Figure Vertex Label]
  \label{def:physics:phyx-mini-0566:phyxminiproblems-problemphyxmini0566-figurevertexlabel}
  \lean{PhyXMiniProblems.ProblemPhyXMini0566.FigureVertexLabel}
  The six vertex labels printed in the supplied figure
\end{definition}
\begin{proof}
  This is the declared model definition of Figure Vertex Label.
\end{proof}
\begin{definition}[Coordinate Axis Label]
  \label{def:physics:phyx-mini-0566:phyxminiproblems-problemphyxmini0566-coordinateaxislabel}
  \lean{PhyXMiniProblems.ProblemPhyXMini0566.CoordinateAxisLabel}
  Coordinate-axis labels visible in the two panels of the figure
\end{definition}
\begin{proof}
  This is the declared model definition of Coordinate Axis Label.
\end{proof}
\begin{definition}[Horizontal Direction]
  \label{def:physics:phyx-mini-0566:phyxminiproblems-problemphyxmini0566-horizontaldirection}
  \lean{PhyXMiniProblems.ProblemPhyXMini0566.HorizontalDirection}
  Direction of the relative-velocity arrow along the common horizontal axis
\end{definition}
\begin{proof}
  This is the declared model definition of Horizontal Direction.
\end{proof}
\begin{definition}[Figure Feature]
  \label{def:physics:phyx-mini-0566:phyxminiproblems-problemphyxmini0566-figurefeature}
  \lean{PhyXMiniProblems.ProblemPhyXMini0566.FigureFeature}
  Named qualitative features visible in the primary raster 566.png
\end{definition}
\begin{proof}
  This is the declared model definition of Figure Feature.
\end{proof}
\begin{definition}[Planar Point Meters]
  \label{def:physics:phyx-mini-0566:phyxminiproblems-problemphyxmini0566-planarpointmeters}
  \lean{PhyXMiniProblems.ProblemPhyXMini0566.PlanarPointMeters}
  A signed Cartesian coordinate readout in metres. This is deliberately a scalar readout of a point in a selected frame, not a replacement for the physical lengths, areas, or speed in the setup
\end{definition}
\begin{proof}
  This is the declared model definition of Planar Point Meters.
\end{proof}
\begin{definition}[squared Distance Meters]
  \label{def:physics:phyx-mini-0566:phyxminiproblems-problemphyxmini0566-squareddistancemeters}
  \lean{PhyXMiniProblems.ProblemPhyXMini0566.squaredDistanceMeters}
  \uses{def:physics:phyx-mini-0566:phyxminiproblems-problemphyxmini0566-planarpointmeters}
  Squared Euclidean separation of two coordinate readouts, in square metres
\end{definition}
\begin{proof}
  This is the declared model definition of squared Distance Meters.
\end{proof}
\begin{definition}[midpoint Meters]
  \label{def:physics:phyx-mini-0566:phyxminiproblems-problemphyxmini0566-midpointmeters}
  \lean{PhyXMiniProblems.ProblemPhyXMini0566.midpointMeters}
  \uses{def:physics:phyx-mini-0566:phyxminiproblems-problemphyxmini0566-planarpointmeters}
  Coordinate midpoint of two planar point readouts
\end{definition}
\begin{proof}
  This is the declared model definition of midpoint Meters.
\end{proof}
\begin{definition}[coordinate Triangle Area Square Meters]
  \label{def:physics:phyx-mini-0566:phyxminiproblems-problemphyxmini0566-coordinatetriangleareasquaremeters}
  \lean{PhyXMiniProblems.ProblemPhyXMini0566.coordinateTriangleAreaSquareMeters}
  \uses{def:physics:phyx-mini-0566:phyxminiproblems-problemphyxmini0566-planarpointmeters}
  Unsigned triangle area obtained from three Cartesian coordinate readouts. The determinant is the oriented double area, so the result is in square metres
\end{definition}
\begin{proof}
  This is the declared model definition of coordinate Triangle Area Square Meters.
\end{proof}
\begin{definition}[Relativistic Triangle Figure]
  \label{def:physics:phyx-mini-0566:phyxminiproblems-problemphyxmini0566-relativistictrianglefigure}
  \lean{PhyXMiniProblems.ProblemPhyXMini0566.RelativisticTriangleFigure}
  \uses{def:physics:phyx-mini-0566:phyxminiproblems-problemphyxmini0566-trianglevertex, def:physics:phyx-mini-0566:phyxminiproblems-problemphyxmini0566-figurevertexlabel, def:physics:phyx-mini-0566:phyxminiproblems-problemphyxmini0566-coordinateaxislabel, def:physics:phyx-mini-0566:phyxminiproblems-problemphyxmini0566-horizontaldirection, def:physics:phyx-mini-0566:phyxminiproblems-problemphyxmini0566-figurefeature}
  Qualitative transcription of the supplied image. The image carries symbolic labels and directions; the numerical 3 m², 15°, and 0.6 c data come from the prose and are linked separately below
\end{definition}
\begin{proof}
  This is the declared model definition of Relativistic Triangle Figure.
\end{proof}
\begin{definition}[Relativistic Triangle Setup]
  \label{def:physics:phyx-mini-0566:phyxminiproblems-problemphyxmini0566-relativistictrianglesetup}
  \lean{PhyXMiniProblems.ProblemPhyXMini0566.RelativisticTriangleSetup}
  \uses{def:physics:phyx-mini-0566:phyxminiproblems-problemphyxmini0566-lengthquantity, def:physics:phyx-mini-0566:phyxminiproblems-problemphyxmini0566-areaquantity, def:physics:phyx-mini-0566:phyxminiproblems-problemphyxmini0566-observerlabel, def:physics:phyx-mini-0566:phyxminiproblems-problemphyxmini0566-trianglevertex, def:physics:phyx-mini-0566:phyxminiproblems-problemphyxmini0566-horizontaldirection, def:physics:phyx-mini-0566:phyxminiproblems-problemphyxmini0566-planarpointmeters, def:physics:phyx-mini-0566:phyxminiproblems-problemphyxmini0566-relativistictrianglefigure}
  The triangle and its independent observables in the two observer frames. The measured area for OPrime is an independent physical quantity. It is not defined from a displayed answer or from a Lorentz-contraction formula. The coordinate and sampling-time fields give the operational data from which each observer's simultaneous area measurement can be related to that physical area
\end{definition}
\begin{proof}
  This is the declared model definition of Relativistic Triangle Setup.
\end{proof}
\begin{definition}[speed Fraction Of Light]
  \label{def:physics:phyx-mini-0566:phyxminiproblems-problemphyxmini0566-speedfractionoflight}
  \lean{PhyXMiniProblems.ProblemPhyXMini0566.speedFractionOfLight}
  \uses{def:physics:phyx-mini-0566:phyxminiproblems-problemphyxmini0566-speedinmeterspersecond, def:physics:phyx-mini-0566:phyxminiproblems-problemphyxmini0566-vacuumspeedoflightinmeterspersecond, def:physics:phyx-mini-0566:phyxminiproblems-problemphyxmini0566-relativistictrianglesetup}
  The dimensionless speed parameter β = v/c
\end{definition}
\begin{proof}
  This is the declared model definition of speed Fraction Of Light.
\end{proof}
\begin{definition}[lorentz Factor]
  \label{def:physics:phyx-mini-0566:phyxminiproblems-problemphyxmini0566-lorentzfactor}
  \lean{PhyXMiniProblems.ProblemPhyXMini0566.lorentzFactor}
  \uses{def:physics:phyx-mini-0566:phyxminiproblems-problemphyxmini0566-relativistictrianglesetup, def:physics:phyx-mini-0566:phyxminiproblems-problemphyxmini0566-speedfractionoflight}
  Physlib's Lorentz factor for the relative speed
\end{definition}
\begin{proof}
  This is the declared model definition of lorentz Factor.
\end{proof}
\begin{definition}[coordinate Area Square Meters]
  \label{def:physics:phyx-mini-0566:phyxminiproblems-problemphyxmini0566-coordinateareasquaremeters}
  \lean{PhyXMiniProblems.ProblemPhyXMini0566.coordinateAreaSquareMeters}
  \uses{def:physics:phyx-mini-0566:phyxminiproblems-problemphyxmini0566-observerlabel, def:physics:phyx-mini-0566:phyxminiproblems-problemphyxmini0566-coordinatetriangleareasquaremeters, def:physics:phyx-mini-0566:phyxminiproblems-problemphyxmini0566-relativistictrianglesetup}
  Coordinate area of the three labelled vertices in one observer frame
\end{definition}
\begin{proof}
  This is the declared model definition of coordinate Area Square Meters.
\end{proof}
\begin{definition}[Matches Relativistic Triangle Scenario]
  \label{def:physics:phyx-mini-0566:phyxminiproblems-problemphyxmini0566-matchesrelativistictrianglescenario}
  \lean{PhyXMiniProblems.ProblemPhyXMini0566.MatchesRelativisticTriangleScenario}
  \uses{def:physics:phyx-mini-0566:phyxminiproblems-problemphyxmini0566-relativistictrianglesetup}
  Observer roles and the rightward relative motion stated in the prose
\end{definition}
\begin{proof}
  This is the declared model definition of Matches Relativistic Triangle Scenario.
\end{proof}
\begin{definition}[Matches Supplied Triangle Figure]
  \label{def:physics:phyx-mini-0566:phyxminiproblems-problemphyxmini0566-matchessuppliedtrianglefigure}
  \lean{PhyXMiniProblems.ProblemPhyXMini0566.MatchesSuppliedTriangleFigure}
  \uses{def:physics:phyx-mini-0566:phyxminiproblems-problemphyxmini0566-relativistictrianglesetup}
  Direct qualitative evidence from phyx\_data/test\_image/566.png: both triangles, the primed and unprimed labels, both coordinate systems, the x ≡ x' mark, the rightward velocity arrow, and the symmetry-axis angle from C
\end{definition}
\begin{proof}
  This is the declared model definition of Matches Supplied Triangle Figure.
\end{proof}
\begin{definition}[Matches Problem Data]
  \label{def:physics:phyx-mini-0566:phyxminiproblems-problemphyxmini0566-matchesproblemdata}
  \lean{PhyXMiniProblems.ProblemPhyXMini0566.MatchesProblemData}
  \uses{def:physics:phyx-mini-0566:phyxminiproblems-problemphyxmini0566-areainsquaremeters, def:physics:phyx-mini-0566:phyxminiproblems-problemphyxmini0566-degreestoradians, def:physics:phyx-mini-0566:phyxminiproblems-problemphyxmini0566-relativistictrianglesetup, def:physics:phyx-mini-0566:phyxminiproblems-problemphyxmini0566-speedfractionoflight}
  The numerical readouts supplied by the problem statement. No value for the area measured by OPrime occurs here
\end{definition}
\begin{proof}
  This is the declared model definition of Matches Problem Data.
\end{proof}
\begin{definition}[Satisfies Rest Equilateral Geometry]
  \label{def:physics:phyx-mini-0566:phyxminiproblems-problemphyxmini0566-satisfiesrestequilateralgeometry}
  \lean{PhyXMiniProblems.ProblemPhyXMini0566.SatisfiesRestEquilateralGeometry}
  \uses{def:physics:phyx-mini-0566:phyxminiproblems-problemphyxmini0566-lengthinmeters, def:physics:phyx-mini-0566:phyxminiproblems-problemphyxmini0566-squareddistancemeters, def:physics:phyx-mini-0566:phyxminiproblems-problemphyxmini0566-midpointmeters, def:physics:phyx-mini-0566:phyxminiproblems-problemphyxmini0566-relativistictrianglesetup}
  The rest-frame coordinate triangle is equilateral, and its symmetry axis runs from C to the midpoint of AB with the stated sine/cosine projections. These are ordinary Euclidean geometry constraints, not an area answer
\end{definition}
\begin{proof}
  This is the declared model definition of Satisfies Rest Equilateral Geometry.
\end{proof}
\begin{definition}[Has Physical Relativistic Parameters]
  \label{def:physics:phyx-mini-0566:phyxminiproblems-problemphyxmini0566-hasphysicalrelativisticparameters}
  \lean{PhyXMiniProblems.ProblemPhyXMini0566.HasPhysicalRelativisticParameters}
  \uses{def:physics:phyx-mini-0566:phyxminiproblems-problemphyxmini0566-lengthinmeters, def:physics:phyx-mini-0566:phyxminiproblems-problemphyxmini0566-areainsquaremeters, def:physics:phyx-mini-0566:phyxminiproblems-problemphyxmini0566-relativistictrianglesetup, def:physics:phyx-mini-0566:phyxminiproblems-problemphyxmini0566-speedfractionoflight}
  Positivity, angle range, and subluminal conditions for the physical branch
\end{definition}
\begin{proof}
  This is the declared model definition of Has Physical Relativistic Parameters.
\end{proof}
\begin{definition}[Satisfies Simultaneous Coordinate Area Measurement]
  \label{def:physics:phyx-mini-0566:phyxminiproblems-problemphyxmini0566-satisfiessimultaneouscoordinateareameasurement}
  \lean{PhyXMiniProblems.ProblemPhyXMini0566.SatisfiesSimultaneousCoordinateAreaMeasurement}
  \uses{def:physics:phyx-mini-0566:phyxminiproblems-problemphyxmini0566-areainsquaremeters, def:physics:phyx-mini-0566:phyxminiproblems-problemphyxmini0566-observerlabel, def:physics:phyx-mini-0566:phyxminiproblems-problemphyxmini0566-trianglevertex, def:physics:phyx-mini-0566:phyxminiproblems-problemphyxmini0566-relativistictrianglesetup, def:physics:phyx-mini-0566:phyxminiproblems-problemphyxmini0566-coordinateareasquaremeters}
  Each observer samples the three vertices simultaneously in that observer's frame, and the physical area readout agrees with the determinant area of those simultaneous coordinates
\end{definition}
\begin{proof}
  This is the declared model definition of Satisfies Simultaneous Coordinate Area Measurement.
\end{proof}
\begin{definition}[Obeys Lorentz Coordinate Contraction]
  \label{def:physics:phyx-mini-0566:phyxminiproblems-problemphyxmini0566-obeyslorentzcoordinatecontraction}
  \lean{PhyXMiniProblems.ProblemPhyXMini0566.ObeysLorentzCoordinateContraction}
  \uses{def:physics:phyx-mini-0566:phyxminiproblems-problemphyxmini0566-trianglevertex, def:physics:phyx-mini-0566:phyxminiproblems-problemphyxmini0566-relativistictrianglesetup, def:physics:phyx-mini-0566:phyxminiproblems-problemphyxmini0566-lorentzfactor}
  Special-relativistic spatial contraction on the simultaneous moving-frame slice. Relative to vertex C, the coordinate parallel to the boost is divided by Physlib's γ(β); the transverse coordinate is invariant. This is a generic governing law. It contains neither a numerical moving area nor an answer-choice label
\end{definition}
\begin{proof}
  This is the declared model definition of Obeys Lorentz Coordinate Contraction.
\end{proof}
\begin{lemma}[area scales by inverse lorentz factor]
  \label{lem:physics:phyx-mini-0566:phyxminiproblems-problemphyxmini0566-area-scales-by-inverse-lorentz-factor}
  \lean{PhyXMiniProblems.ProblemPhyXMini0566.area_scales_by_inverse_lorentz_factor}
  \uses{def:physics:phyx-mini-0566:phyxminiproblems-problemphyxmini0566-areainsquaremeters, def:physics:phyx-mini-0566:phyxminiproblems-problemphyxmini0566-relativistictrianglesetup, def:physics:phyx-mini-0566:phyxminiproblems-problemphyxmini0566-lorentzfactor, def:physics:phyx-mini-0566:phyxminiproblems-problemphyxmini0566-satisfiessimultaneouscoordinateareameasurement, def:physics:phyx-mini-0566:phyxminiproblems-problemphyxmini0566-obeyslorentzcoordinatecontraction}
  The determinant area law and coordinate contraction imply the generic standard-model scaling A' = A / γ(β). This is a derived conclusion rather than a field of either governing-law structure
\end{lemma}
\begin{proof}
  The conclusion follows from the governing relations and figure readouts named in the statement of area scales by inverse lorentz factor.
\end{proof}
\begin{lemma}[standard model area is two point four]
  \label{lem:physics:phyx-mini-0566:phyxminiproblems-problemphyxmini0566-standard-model-area-is-two-point-four}
  \lean{PhyXMiniProblems.ProblemPhyXMini0566.standard_model_area_is_two_point_four}
  \uses{def:physics:phyx-mini-0566:phyxminiproblems-problemphyxmini0566-areainsquaremeters, def:physics:phyx-mini-0566:phyxminiproblems-problemphyxmini0566-relativistictrianglesetup, def:physics:phyx-mini-0566:phyxminiproblems-problemphyxmini0566-matchesproblemdata, def:physics:phyx-mini-0566:phyxminiproblems-problemphyxmini0566-satisfiessimultaneouscoordinateareameasurement, def:physics:phyx-mini-0566:phyxminiproblems-problemphyxmini0566-obeyslorentzcoordinatecontraction}
  For the stated 3 m² rest area and β = 3/5, the standard simultaneous Lorentz-contraction model gives A' = 12/5 m² = 2.4 m²
\end{lemma}
\begin{proof}
  The conclusion follows from the governing relations and figure readouts named in the statement of standard model area is two point four.
\end{proof}
\begin{definition}[Answer Choice]
  \label{def:physics:phyx-mini-0566:phyxminiproblems-problemphyxmini0566-answerchoice}
  \lean{PhyXMiniProblems.ProblemPhyXMini0566.AnswerChoice}
  Labels of the four square-metre-valued choices in the source
\end{definition}
\begin{proof}
  This is the declared model definition of Answer Choice.
\end{proof}
\begin{definition}[displayed Area Square Meters]
  \label{def:physics:phyx-mini-0566:phyxminiproblems-problemphyxmini0566-displayedareasquaremeters}
  \lean{PhyXMiniProblems.ProblemPhyXMini0566.displayedAreaSquareMeters}
  \uses{def:physics:phyx-mini-0566:phyxminiproblems-problemphyxmini0566-answerchoice}
  Area in square metres printed beside each answer label
\end{definition}
\begin{proof}
  This is the declared model definition of displayed Area Square Meters.
\end{proof}
\begin{definition}[recorded Dataset Answer]
  \label{def:physics:phyx-mini-0566:phyxminiproblems-problemphyxmini0566-recordeddatasetanswer}
  \lean{PhyXMiniProblems.ProblemPhyXMini0566.recordedDatasetAnswer}
  \uses{def:physics:phyx-mini-0566:phyxminiproblems-problemphyxmini0566-answerchoice}
  The answer label recorded in the supplied dataset
\end{definition}
\begin{proof}
  This is the declared model definition of recorded Dataset Answer.
\end{proof}
\begin{definition}[Matches Displayed Area Choice]
  \label{def:physics:phyx-mini-0566:phyxminiproblems-problemphyxmini0566-matchesdisplayedareachoice}
  \lean{PhyXMiniProblems.ProblemPhyXMini0566.MatchesDisplayedAreaChoice}
  \uses{def:physics:phyx-mini-0566:phyxminiproblems-problemphyxmini0566-areainsquaremeters, def:physics:phyx-mini-0566:phyxminiproblems-problemphyxmini0566-relativistictrianglesetup, def:physics:phyx-mini-0566:phyxminiproblems-problemphyxmini0566-answerchoice, def:physics:phyx-mini-0566:phyxminiproblems-problemphyxmini0566-displayedareasquaremeters}
  An answer choice agrees with the independently represented moving-frame area. Unfolding this definition leaves a substantive equality to the measured area; it does not make any choice correct by definition
\end{definition}
\begin{proof}
  This is the declared model definition of Matches Displayed Area Choice.
\end{proof}
\begin{theorem}[recorded answer disagrees with standard model]
  \label{thm:physics:phyx-mini-0566:phyxminiproblems-problemphyxmini0566-recorded-answer-disagrees-with-standard-model}
  \lean{PhyXMiniProblems.ProblemPhyXMini0566.recorded_answer_disagrees_with_standard_model}
  \uses{def:physics:phyx-mini-0566:phyxminiproblems-problemphyxmini0566-relativistictrianglesetup, def:physics:phyx-mini-0566:phyxminiproblems-problemphyxmini0566-matchesproblemdata, def:physics:phyx-mini-0566:phyxminiproblems-problemphyxmini0566-satisfiessimultaneouscoordinateareameasurement, def:physics:phyx-mini-0566:phyxminiproblems-problemphyxmini0566-obeyslorentzcoordinatecontraction, def:physics:phyx-mini-0566:phyxminiproblems-problemphyxmini0566-recordeddatasetanswer, def:physics:phyx-mini-0566:phyxminiproblems-problemphyxmini0566-matchesdisplayedareachoice}
  The source metadata records choice B, whose displayed value is 7/5 m². Under the same stated data and standard contraction law, that choice does not match the independently represented moving-frame area. This conclusion makes the source discrepancy machine-visible without treating either the recorded label or its displayed value as physical input data
\end{theorem}
\begin{proof}
  The conclusion follows from the governing relations and figure readouts named in the statement of recorded answer disagrees with standard model.
\end{proof}
% --- Archon declaration topology end ---
% --- Archon physics formalization source end ---
